\documentclass[11pt,a4paper]{article}
\usepackage{amsmath,amsthm}
\usepackage[utf8]{inputenc}
\newtheorem{thm}{Theorem}[section]
\newtheorem{lem}[thm]{Lemma}
\newtheorem*{rem}{Remark}
\newcommand{\R}{\mathbb{R}}
\newcommand{\norm}[1]{\lVert #1 \rVert}
\DeclareMathOperator{\tr}{tr}
\def\half{\frac{1}{2}}
\newenvironment{note}[1][Note]{\textbf{#1.}}{}

\begin{document}
\section{Introduction}\label{sec:intro}
Text\footnote{A note.\label{fn:one}} and more\footnote{Two.}.
\begin{equation}\label{eq:one}
  a = b
\end{equation}
\begin{align}
  x &= 1 \label{eq:x}\\
  y &= 2 \nonumber\\
  z &= 3 \tag{star}\label{eq:z}\\
  w &= 4 \label{eq:w}\\
\end{align}
\begin{equation*}
  c = d
\end{equation*}
\begin{subequations}\label{eq:sub}
\begin{align}
  p &= 1 \label{eq:p}\\
  q &= 2 \label{eq:q}
\end{align}
\end{subequations}
\begin{equation}
\begin{split}
  e &= f\\
    &= g
\end{split}\label{eq:split}
\end{equation}
\subsection{Background}\label{sec:bg}
\begin{thm}[Pythagoras]\label{thm:py}
  $a^2+b^2=c^2$.
\end{thm}
\begin{lem}\label{lem:a}
  Easy.
\end{lem}
\begin{rem}
  Unnumbered.
\end{rem}
\subsubsection{Deep}\label{sec:deep}
\paragraph{Para}\label{sec:para}
Text.
\section*{Unnumbered}
\section{Figures}\label{sec:fig}
\begin{figure}[h]
  \centering
  \caption[Short]{A figure.}\label{fig:a}
\end{figure}
\begin{table}
  \caption{A table.}\label{tab:a}
\end{table}
\begin{figure*}
  \caption{Wide.}\label{fig:b}
\end{figure*}
\begin{thm}\label{thm:two}
  Another.
\end{thm}
See \ref{sec:intro}, \eqref{eq:x} and \cite[p.~3]{knuth,lamport}.
\setcounter{secnumdepth}{1}
\subsection{Not numbered now}\label{sec:nonum}
\appendix
\section{Proofs}\label{sec:proofs}
\subsection{First}\label{sec:proofs1}
\begin{equation}\label{eq:app}
  u = v
\end{equation}
\bibliographystyle{plain}
\bibliography{refs,more}
\end{document}
